% Generated draft: requires mathematical and source review.
\documentclass[11pt]{article}
\usepackage[margin=25mm]{geometry}
\usepackage{fontspec}
\setmainfont{Latin Modern Roman}
\usepackage{amsmath,amssymb,mathtools,unicode-math}
\setmathfont{Latin Modern Math}
\usepackage{enumitem,needspace,xurl,hyperref}
\setlength{\parindent}{0pt}
\setlength{\parskip}{5pt}
\setlength{\emergencystretch}{3em}
\newcommand{\N}{\mathbb{N}}
\newcommand{\Z}{\mathbb{Z}}
\newcommand{\Q}{\mathbb{Q}}
\newcommand{\R}{\mathbb{R}}
\newcommand{\C}{\mathbb{C}}
\newcommand{\F}{\mathbb{F}}
\newcommand{\A}{\mathbb{A}}
\newcommand{\PP}{\mathbb{P}}
\newcommand{\E}{\mathbb{E}}
\newcommand{\eps}{\varepsilon}
\newcommand{\dd}{\,\mathrm{d}}
\begin{document}
\section*{UnsolvedMath ID: 20003176}
\textbf{Source problem code:} catalogue entry
\subsection{Definitions and mathematical statement}
A connected manifold $M$ is a $K(\pi,1)$, or aspherical, if $\pi=\pi_1(M)$ and $\pi_j(M)=0$ for every $j\geq2$; for manifolds this is equivalent to contractibility of the universal cover. Problem 51 asks whether a generic manifold is aspherical, explicitly leaving generic to be determined. Thus no truth value is fixed without a class of manifolds and a topology, probability distribution, enumeration or other genericity model. The workshop's curvature theme does not by itself specify that every proposed model is restricted to nonnegative curvature.

The aid records several model distinctions. Changing the Riemannian metric on a fixed smooth manifold cannot change asphericity, so Baire-generic metrics there give either the whole space or the empty locus. In a natural topological disjoint union across diffeomorphism types, the union of aspherical components is clopen; when both kinds occur in a Baire model, it is a proper nondense locus rather than a residual one. For each fixed dimension at least three, the aid also reports that enumeration densities across unrestricted manifold types can be calibrated arbitrarily by choosing an enumeration, using infinite collections of aspherical and nonaspherical types. This is a dependence on the ordering, not an intrinsic probability or a curvature-restricted assertion.

A concrete affirmative probabilistic model in the aid is the Pippenger--Schleich quotient construction. Take an even number $n$ of oriented triangles, choose a uniformly random perfect matching of their $3n$ sides, and glue matched sides by the orientation-reversing affine identifications appropriate to an oriented surface. The resulting quotient is allowed to be disconnected and need not be an abstract simplicial complex without repeated identifications. The exact quotient conventions must match the cited model. The reported estimate, as $n\to\infty$ through even integers, is
\[
 \Pr\{\text{the quotient is connected and a }K(\pi,1)\}
 =1-O\!\left(\frac{\log n}{n}\right).
\]
This statement is specifically about closed oriented surface quotients, not all dimensions or every random-triangulation distribution.

The aid attributes the input to Pippenger--Schleich's expected vertex count and connectedness theorem. If $V_n$ is the quotient-vertex count, it reports $\E V_n=O(\log n)$. On the connected event, Euler characteristic is $V_n-n/2$. A sphere therefore requires $V_n=n/2+2$, and Markov's inequality gives
\[
 \Pr\{\text{connected sphere}\}
 \leq\frac{\E V_n}{n/2+2}
 =O\!\left(\frac{\log n}{n}\right).
\]
Combining with the cited connectedness control gives the displayed bound, since a connected closed orientable surface of genus at least one is aspherical. The exact theorem constants and connectedness estimate are not reproduced in the record. This is a reported model-specific result, not a resolution of an undefined model-independent genericity claim.
\subsection{Short English statement}
Genericity must be chosen before the question is meaningful. The aid gives a precise random oriented-surface model in which connected aspherical quotients occur with the stated high probability. Fixed-manifold metric genericity and arbitrary enumeration density address different questions.
\subsection{Sources}
\begin{itemize}
\item Local UnsolvedMath registry, permanent ID 20003176, source code catalogue entry; \texttt{aim-topology-notes.json}, notes[263].
\item AIM, \emph{Manifolds with nonnegative sectional curvature}, page 4, ``8. Miscellaneous Problems,'' Problem 51: \url{https://aimath.org/WWN/nnsectcurvature/nnsectcurvature.pdf}. The supplied verification repairs only the line-break hyphen in ``determined.''
\item Pippenger--Schleich uniform triangle-quotient model, expected vertex count and connectedness theorem, named by the research summary without full bibliography or exact theorem constants.
\item Machine-generated research attempt 1: partial result, original status unknown at the run, rigor labelled ``proved,'' low novelty confidence; registry partially solved and corrected verified statement. Not peer reviewed; independently verify before citation. Registry updated 14 August 2026.
\end{itemize}
\end{document}
